% Generated by paper-covert from output/manuscript/sections_revised/03_related_work.md

\section{Related Work}\label{related-work}

\subsection{Global crop-mapping infrastructure and foundation models}\label{global-crop-mapping-infrastructure-and-foundation-models}

A growing ecosystem of global and continental crop-mapping products now provides open, regularly updated coverage. WorldCereal \citep{vantricht2023worldcereal} is the most directly comparable: a dynamic, open-source system producing seasonal maize, winter-cereal, and spring-cereal classifications and a general temporary-cropland extent layer from Sentinel-1 and Sentinel-2, tiled across 106 globally defined agro-ecological zones. It has no Canola or Legume class, reflecting its design priorities toward globally widespread cereal and maize systems rather than the Southern Hemisphere oilseed-and-pulse rotations that structure the Australian grains sector. Foundation-model embeddings have also been proposed as a general-purpose alternative to hand-engineered spectral features: Presto \citep{tseng2024presto} pretrains a lightweight transformer on multi-modal pixel time series (optical, radar, climate, terrain) for downstream fine-tuning, and AlphaEarth \citep{brown2025alphaearth} extends this idea to a general-purpose annual embedding at global scale. Generic land-use-land-cover products such as Esri/Impact Observatory's global layer \citep{karra2021esri} classify cropland as a single category without species or crop-type resolution at all.

This paper relates to this cluster as a local adaptation-and-comparison study with two genuine negative results rather than a novel infrastructure contribution. We fine-tuned Presto on our own field-verified labels and found it lost to hand-built, day-of-year-binned spectral indices (NDVI, a canola-specific flowering index, and a chlorophyll-fluorescence-related index) at this label volume and class granularity --- macro F1 0.775 versus 0.821 for the hand-built feature set, and markedly worse on canola detection specifically (74.2\% versus 89.0\% at a 5\% false-positive rate). Combining Presto embeddings with the hand-built features gave no improvement over the hand-built features alone (0.824 versus 0.821, within the noise of a size-matched control comparison). The most plausible cause is not that Presto's representation is unsuited to this problem, but that three of its nine expected input channel groups --- Sentinel-1 radar, ERA5 climate, and SRTM terrain --- are absent from our input data and must be passed as masked rather than genuinely observed; this is a data-availability limitation of our deployment, not necessarily a finding about Presto's architecture, and we report it as such. We also directly compared our national map against WorldCereal's 2021 products at the pixel level (Section 5.5): because WorldCereal has no Canola or Legume class, only the presence/absence extent layer and the wintercereals layer are checkable against our output, and both comparisons are reported with their full caveats rather than treated as ground truth.

\subsection{Australia-specific crop-type classification precedents}\label{australia-specific-crop-type-classification-precedents}

Two studies are the closest direct precedents. \citet{sharma2026wa} built an in-season Sentinel-2 classifier for six crop classes in Western Australia, using day-after-sowing-derived pseudo-labels for the bulk of training data supplemented by a smaller field-verified test set. \citet{alshammari2024mdb} fused Sentinel-1, Sentinel-2, and MODIS to distinguish cereals from canola across the Murray-Darling Basin, training on harvester-derived yield-map labels rather than recorded sowings. Both reach classification accuracies above 90\% within their respective study regions, and both represent substantial prior engineering effort. Neither, however, is continent-wide, and neither trains exclusively on field-verified ground truth: \citeauthor{sharma2026wa}'s pseudo-labels are model-derived, and \citeauthor{alshammari2024mdb}'s harvester labels are opportunistic records of what a combine harvested, not confirmed sowings. This paper's methodological position relative to both is to trade breadth of species resolution --- three broad groups rather than six specific crops --- for two things neither precedent has together: continent-wide coverage and training exclusively on GRDC National Variety Trial records, which are field-recorded sowings rather than another model's output or an opportunistic byproduct of harvest operations.

A third, operationally oriented precedent is CSIRO's Graincast system \citep{lawes2022graincast}, which combines a field-boundary product (ePaddocks) with a satellite-and-climate-driven yield model to monitor crop production across the Australian grainbelt at national scale. Graincast's boundary-plus-yield pattern is the closest operational analogue to what this paper attempts, and the original project brief for this work specified ePaddocks as the intended boundary source. In practice, this project replaced ePaddocks with Segment Anything Model (SAM)-based segmentation (Section 2.3) early in development; Graincast is cited here as the closest prior art for the general pattern, not as a method this paper builds on directly.

\subsection{Field and paddock boundary delineation}\label{field-and-paddock-boundary-delineation}

Paddock boundaries were delineated with the Segment Anything Model (SAM; \citealp{kirillov2023sam}), a general-purpose, promptable image-segmentation foundation model, applied here via the SAMGeo/PaddockTS geospatial wrapper with stock parameters. This is a departure from the project's original plan to use ePaddocks (Section 2.2) as the boundary source. Fields of the World \citep{kerner2024fields} is the most directly relevant methodological comparison point among purpose-built agricultural boundary-delineation methods: a cloud-robust Sentinel-1-and-Sentinel-2 field-boundary model evaluated against a multi-country benchmark. Its published benchmark does not include Australia, and its own documentation notes that only a model-predicted global product, not the curated benchmark itself, extends to the Australian continent; this absence of Australian ground-truth boundary data was itself a factor in the decision to use a general-purpose segmentation model rather than attempt to adapt a boundary-specific method without a local evaluation set.

\subsection{Accuracy-assessment standards for crop and land-cover maps}\label{accuracy-assessment-standards-for-crop-and-land-cover-maps}

\citet{olofsson2014goodpractices} set out the field's stated good-practice standard for map accuracy assessment and area estimation: stratified probability sampling of reference data, error matrices, and confidence intervals on both accuracy and estimated area. This standard is named here explicitly because it is \emph{not} what this paper's validation does, and the reason matters for how the results should be read. Olofsson-style area estimation requires a probability sample of reference labels stratified by map class; no such sample of field-verified crop-species labels exists for Australia at the density this would require, because the GRDC/NVT trial network is a purposively sited agronomic trial program, not a probability sample of the landscape. In its place, this paper validates against ABS sown-area statistics aggregated to the SA2 level --- an independent, non-training-derived, but itself uncertain reference (ABS and ABARES, its sister agency, disagree with each other by a median 6.8\% on national sown area) --- reporting composition and area-ratio comparisons rather than a stratified accuracy assessment with confidence intervals. We flag this explicitly as a methodological limitation (Section 7) rather than presenting the ABS comparison as if it satisfied the Olofsson-style standard.

\subsection{Positioning relative to the nearest competing approaches}\label{positioning-relative-to-the-nearest-competing-approaches}

The nearest competing approaches are, respectively, WorldCereal for global coverage without species resolution, and \citet{sharma2026wa} and \citet{alshammari2024mdb} for Australian species-level classification without continent-wide, exclusively field-verified training data. This paper's differentiation along the specific dimension of ground-truth provenance and geographic coverage --- continent-wide, three-class, trained end-to-end on field-recorded sowings --- is a genuine methodological edge purchased at the explicit cost of coarser species resolution than either Australian precedent and of not solving global species generalizability the way a foundation-model approach might eventually. We report this trade-off directly rather than only in retrospect: the original \textasciitilde10-species scope was not reachable at the available field-verified label volume (9-class macro F1 0.38), and the 3-group descoping is a finding about label-volume requirements for field-verified species classification, stated here in the Introduction and Methods rather than discovered by the reader only in the Results.

\emph{Note on data lineage}: \citet{newmanfurbank2021} describes the same underlying GRDC National Variety Trial network this paper's ground truth is drawn from. It is not an independent public dataset, and this paper does not point readers toward it as a substitute or independent replication source; it is cited here only to acknowledge the shared data lineage.
